\section{Metatheory}
\label{sec:metatheory}

This section states the main metatheoretical results for \thecalculus. All
definitions from the previous sections and all theorems of this section are
formalized and proved in Agda. We mark each mechanized theorem with a little
Agda bird after the theorem name. The supplement names the Agda module of each
theorem.

\subsection{Type Safety for Declarative Typing}

Preservation
yields the usual safety invariant across reduction steps.

\begin{theorem}[Expression preservation\;\AgdaBird]
  If\/ $\HasType \SessCtx \E \T$ and $\E \EReduce \E*$, then
  $\HasType \SessCtx {\E*} \T$.
\end{theorem}

\begin{theorem}[Process preservation\;\AgdaBird]
  If\/ $\ProcHasType \SessCtx \Proc$ and $\Proc \PReduce \Proc*$, then
  $\ProcHasType \SessCtx {\Proc*}$.
\end{theorem}

We complement preservation with progress. Expressions are not guaranteed to
reduce to values but may be blocked on a constant applied to a value. The
reduction of constants is part of the process semantics.

\begin{theorem}[Expression progress\;\AgdaBird]
  \label{thm:expression-progress}
  If\/ $\HasType \SessCtx \E \T$, then either $\E$ is a value, or there exist
  an evaluation context $\EC$, constant $\C$ and value $\V$ such that
  $\E=\ECapp{\EApp\C\V}$, or there exists an $\E*$ such that $\E \EReduce \E*$.
\end{theorem}

A process is blocked if every thread is blocked and no restriction can fire a
channel rule. The blocked-process predicate $\Blocked\Proc$ is defined in
\cref{fig:blocked-expr-proc}. Progress holds for closed processes and
guarantees reduction only for processes that are not blocked. Thus, deadlocks
remain possible.

\begin{theorem}[Process progress\;\AgdaBird]
  \label{thm:process-progress}
  If\/ $\ProcHasType \CtxEmpty \Proc$, then $\Proc \Congruent \PExp{\CUnit}$,
  or $\Blocked\Proc$, or there exists a $\Proc*$ such that
  $\Proc \PReduce \Proc*$.
\end{theorem}

\subsection{Simulation for the Translation}

We relate source and target reduction through the image relation
$\Proc \ImageRel \TProcQ$, which holds when the configuration $\TProcQ$
represents the closed, well-typed process $\Proc$ up to the renaming of
primitive channels and the ordering of threads. The translation produces an
image, $\Proc \ImageRel \TransP\Proc$. Both simulations match one source
step with one target step.

\begin{theorem}[Forward simulation\;\AgdaBird]
  If\/ $\ProcHasType \CtxEmpty \Proc$, $\Proc \ImageRel \TProcQ$, and
  $\Proc \PReduce \Proc*$, then there exists $\TProcQ*$ such that
  $\TProcQ \TPReduce \TProcQ*$ and $\Proc* \ImageRel \TProcQ*$.
\end{theorem}

The converse direction holds up to the equivalence $\Bisim$ that identifies
configurations differing only in the numbering of the synchronization slots of
one endpoint. The target rule for a right split may place its new slot at any
position, whereas the source rule \textsc{R-RSplit} fixes that position.

\begin{theorem}[Backward simulation\;\AgdaBird]
  If\/ $\ProcHasType \CtxEmpty \Proc$, $\Proc \ImageRel \TProcQ$,
  $\TProcQ \Bisim \TProcR$, and $\TProcR \TPReduce \TProcR*$, then there
  exist $\Proc*$ and $\TProcQ*$ such that $\Proc \PReduce \Proc*$,
  $\Proc* \ImageRel \TProcQ*$, and $\TProcQ* \Bisim \TProcR*$.
\end{theorem}


\subsection{Soundness and Completeness for Algorithmic Typing}

We now connect algorithmic typing with declarative typing. Soundness says
that any solved algorithmic judgment denotes a valid declarative typing, while
completeness says that declarative typings can be reconstructed algorithmically.

\begin{theorem}[Algorithmic typing soundness\;\AgdaBird]
  If either $\CheckType \E \AT \Eff$ or $\InferType \E \AT \Eff$
  and $\sigma \models \CSet$ is a closing substitution for the unification
  variables of $\Ctx$, $\E$ and $\AT$, then
  $\HasType[\Eff] {\Ctx{[\sigma]}} {\E{[\sigma]}} {\AT{[\sigma]}}$.
\end{theorem}

Completeness starts from a context, an expression and a type without
unification variables. It requires that every variable of linear type occurs
at most once in $\Ctx$. This condition is necessary, because the restriction
of $\Ctx$ to the free variables of a subexpression keeps every occurrence of
a repeated variable, which no algorithmic rule for variables can consume. The
supplement gives a counterexample. Both theorems produce an \emph{annotated
variant} $\E*$ of $\E$, which adds type annotations $\E:\T$ to subterms and
leaves $\E$ otherwise unchanged; the proofs insert them exactly at checking
forms in inference position, with the declarative type of that subterm.

\begin{theorem}[Checking completeness\;\AgdaBird]
  If\/ $\HasType[\Eff] \Ctx \E \T$, then there is an annotated variant
  $\E*$ of $\E$ with $\CheckType {\E*} \T {\Eff*}$ for some
  constraint set $\CSet$, effect $\Eff* \le \Eff$, and closing substitution
  $\sigma \models \CSet$.
\end{theorem}

\begin{theorem}[Inference completeness\;\AgdaBird]
  If\/ $\HasType[\Eff] \Ctx \E \T$, then there is an annotated variant
  $\E*$ of $\E$ with $\InferType {\E*} \AT {\Eff*}$ for some type $\AT$,
  constraint set $\CSet$, effect $\Eff* \le \Eff$, and closing substitution
  $\sigma \models \CSet$ with $\AT{[\sigma]} \simeq \T$.
\end{theorem}

%%% Local Variables:
%%% mode: latex
%%% TeX-master: "../popl2027.tex"
%%% End:
